\documentclass[border=10pt]{standalone}
\usepackage[european]{circuitikz}
\begin{document}
\begin{tikzpicture}[x=1cm, y=1cm, font=\normalsize, line width=0.8pt]
  % 图 8.02：CP 四脉冲。横轴 0–100% 等分 20 格，每格 0.55 cm。
  % 沿的位置（照 078 §3 与课程站 10-26 页转写表）：
  %   上升 10% / 30% / 50% / 70%  -> x = 1.1 / 3.3 / 5.5 / 7.7
  %   下降 25% / 40% / 60% / 80%  -> x = 2.75 / 4.4 / 6.6 / 8.8
  % 总宽 20 格 = 11.0 cm；电平 0 -> y=0，1 -> y=0.75
  \def\seg{0.55}
  \def\hi{0.75}
  \def\total{11.0}

  \node[anchor=east] at (-0.3,\hi/2) {CP};

  % 电平串：0 |1| 0 |1| 0 |1| 0 |1| 0  （段长依次为 2,3,1,2,2,2,2,2,4 格）
  \draw (0,0) -- (2*\seg,0);
  \draw (2*\seg,0) -- (2*\seg,\hi) -- (5*\seg,\hi);
  \draw (5*\seg,\hi) -- (5*\seg,0) -- (6*\seg,0);
  \draw (6*\seg,0) -- (6*\seg,\hi) -- (8*\seg,\hi);
  \draw (8*\seg,\hi) -- (8*\seg,0) -- (10*\seg,0);
  \draw (10*\seg,0) -- (10*\seg,\hi) -- (12*\seg,\hi);
  \draw (12*\seg,\hi) -- (12*\seg,0) -- (14*\seg,0);
  \draw (14*\seg,0) -- (14*\seg,\hi) -- (16*\seg,\hi);
  \draw (16*\seg,\hi) -- (16*\seg,0) -- (\total,0);

  % 沿的位置标注（8 个沿，与转写表一致）
  \node[anchor=south, font=\small] at (2*\seg,\hi+0.12) {$\uparrow_1$};
  \node[anchor=north, font=\small] at (5*\seg,-0.28) {$\downarrow_1$};
  \node[anchor=south, font=\small] at (6*\seg,\hi+0.12) {$\uparrow_2$};
  \node[anchor=north, font=\small] at (8*\seg,-0.28) {$\downarrow_2$};
  \node[anchor=south, font=\small] at (10*\seg,\hi+0.12) {$\uparrow_3$};
  \node[anchor=north, font=\small] at (12*\seg,-0.28) {$\downarrow_3$};
  \node[anchor=south, font=\small] at (14*\seg,\hi+0.12) {$\uparrow_4$};
  \node[anchor=north, font=\small] at (16*\seg,-0.28) {$\downarrow_4$};

\end{tikzpicture}
\end{document}
